% dynamics/optimization_in_state_space-DE.tex
\chapter{Optimierung erzeugt Kandidaten}
\label{chap:optimization-in-state-space}

\section*{Ein Ziel benötigt einen Verantwortlichen}

Die vorhergehenden Kapitel können qualifizierte Modellausgaben und Vergleiche
erzeugen. Sie machen keine Größe von selbst erstrebenswert. Kurze
Übergangslänge, reduziertes Gleichgewichtsresiduum, Komfort, Verschleiß,
Kosten und Energie sind verschiedene Ziele, solange eine verantwortliche
Quelle sie nicht ausdrücklich aggregiert.

\begin{principle}[Prinzip qualifizierter Optimierung]
\label{princ:state-space-optimization}
Optimierung beginnt erst, wenn Problemidentität, Zweck,
Entscheidungsvariablen, feste Eingaben, zulässige Menge, Modell und Ziel
erklärt sind.
\end{principle}

\begin{principle}[Kandidatenprinzip]
\label{princ:trajectory-shaping}
Ein Solver liefert einen Berechnungskandidaten unter dieser Erklärung, keine
Engineering Decision.
\end{principle}

\begin{figure}[htbp]
\centering
\begin{tikzpicture}[
  node distance=5mm and 6mm,
  box/.style={draw,rounded corners,align=center,minimum height=9mm,text width=27mm,font=\scriptsize},
  arrow/.style={->,thick}
]
\node[box] (problem) {Problemidentität\\Zweck und Freiheiten};
\node[box,right=of problem] (model) {identifiziertes Modell \(M\)\\feste Eingaben};
\node[box,right=of model] (objective) {erklärte \(J,\mathcal F\)\\Einheiten und Quelle};
\node[box,below=of objective] (candidate) {Berechnungskandidat\\Werte und Residuen};
\node[box,left=of candidate] (evaluation) {qualifizierte Bewertung\\Evidenz und Regeln};
\node[box,left=of evaluation] (decision) {Engineering Decision\\Autorität und Nachweis};
\draw[arrow] (problem)--(model);
\draw[arrow] (model)--(objective);
\draw[arrow] (objective)--(candidate);
\draw[arrow] (candidate)--(evaluation);
\draw[arrow] (evaluation)--(decision);
\end{tikzpicture}

\caption{Abhängigkeit vom identifizierten Problem und Modell über einen nichtautoritativen Kandidaten und seine Bewertung zu einer möglichen Engineering Decision.}
\label{fig:model-candidate-decision}
\end{figure}

\section{Problemvertrag}

\begin{definition}[Laufzeitzustandstrajektorie]
\label{def:runtime-state-trajectory}
Eine Laufzeitzustandstrajektorie ist eine identifizierte Modellausgabe
\(x_M(t)\), die ein erklärtes Optimierungsproblem verwendet. Sie ist kein
universeller AIM-Zustand.
\end{definition}

\begin{definition}[Laufzeitzustandsentwicklung]
\label{def:optimization-runtime-state-evolution}
Ihre Entwicklung wird durch \(\mathcal D_M\) aus festen und variablen Eingaben
unter angegebenen Anfangsbedingungen und Anwendbarkeit erzeugt.
\end{definition}

Erst dann ist die bekannte Form sinnvoll:
\[
\min_{x\in\mathcal F}J(x).
\]
Dabei bezeichnet \(x\) den erklärten Entscheidungsvektor, keinen
unspezifizierten Zustand. Der Problemnachweis enthält
\[
\begin{aligned}
\mathcal P=(\, &id,\ purpose,\ x,\ b,\ f,\ d,\ M,\ \mathcal F,\ J,\\
              &units,\ scaling,\ source,\ applicability,\ provenance\,),
\end{aligned}
\]
wobei \(b\) Grenzen, \(f\) feste Eingaben und \(d\) abgeleitete Größen sind.
Offene Freiheiten bleiben sichtbar, statt stillschweigend belegt zu werden.

\section{Freiheitsgrade}

Für eine beschränkte Kontinuitätsänderung lautet ein vertretbares Beispiel
\[
x=(L,\alpha),\qquad
\mathcal F=\{x\mid c_{\mathrm{continuity}}(x)=0,
L_{\min}\le L\le L_{\max},
\alpha_{\min}\le\alpha\le\alpha_{\max}\}.
\]
Bekannte Endzustände und gewählte Übergangsfamilie sind fest; \(L,\alpha\)
sind innerhalb ihrer Grenzen frei; Koeffizienten und End-Pose werden
abgeleitet; Kontinuitätsresiduen sind beschränkt. Überhöhung,
Fahrzeugparameter und operatives Geschwindigkeitsprofil bleiben offen, sofern
sie nicht bewusst ergänzt werden.

Der aktuelle ufAIM/AXTRAN-Kontinuitätssolver ist Implementierungsevidenz für
eine solche beschränkte Kontinuitätslösung. Er belegt keine allgemeine
Fahrzeug-, Überhöhungs-, Komfort- oder Netzoptimierung und keine automatische
ingenieurtechnische Auswahl.

\section{Ziel und Skalierung}

\begin{definition}[Laufzeitzustandsfunktional]
\label{def:runtime-state-functional}
Ein Laufzeitzustandsfunktional ist ein erklärtes Ziel aus einer
identifizierten Modellausgabe, dessen Einheiten, Skalierung, Quelle und Zweck
dokumentiert sind.
\end{definition}

\begin{definition}[Operative Trajektorienformung]
\label{def:optimization-trajectory-shaping}
Operative Trajektorienformung ist eine mögliche Optimierungsformulierung,
keine kanonische AIM-Interpretation.
\end{definition}

Die beiden Ziele \(J_\kappa\) und \(J_r\) in
\cref{chap:vienna-reinterpretation} zeigen, warum Skalarwerte nicht
zweckübergreifend vergleichbar sind. Eine gewichtete Aggregation
\[
J=w_\kappa J_\kappa/J_{\kappa,0}+w_rJ_r/J_{r,0}
\]
ist nur sinnvoll, wenn Referenzskalen, Gewichte und ihre verantwortliche
Quelle dokumentiert sind. Der Solver besitzt sie nicht.

\section{Kandidatenvertrag}

Ein Berechnungskandidat enthält
\[
\begin{aligned}
C=(\, &\mathcal P.id,\ M.id,\ x,\ b,\ \mathcal F,\ J,\ x^\star,\\
      &residuals,\ convergence,\ provenance,\ applicabilityStatus\,).
\end{aligned}
\]
Bei Rückgabe des Solvers ist die Anwendbarkeit noch unbewertet und der
Kandidat ausdrücklich nichtautoritativ.

\[
\boxed{\text{berechnet}
\neq\text{anwendbar}
\neq\text{angenommen}
\neq\text{autoritativ}}
\]

Mehrere Kandidaten können mathematisch zulässig, unter verschiedenen Zielen
optimal, ohne Zweck unvergleichbar, durch spätere Evidenz ungültig oder nach
einer Alignment- oder Realisierungsänderung veraltet sein. Eine skalare
Kennzahl liefert keine universelle ingenieurtechnische Ordnung.

\begin{definition}[Dünn besetzte Laufzeitoptimierung]
\label{def:sparse-runtime-optimization}
Dünn besetzte Laufzeitoptimierung ist lediglich eine mögliche
Berechnungstechnik, die eine erklärte dünn besetzte Abhängigkeitsstruktur
nutzt; eine allgemeine Leistungs- oder Fähigkeitsaussage wird nicht erhoben.
\end{definition}

\begin{definition}[Operativ zulässige Trajektorie]
\label{def:operationally-admissible-trajectory}
Eine operativ zulässige Trajektorie ist ein Bewertungsergebnis unter
identifizierten Regeln, Zweck, Evidenz und Anwendbarkeit. Zulässigkeit in
\(\mathcal F\) allein begründet sie nicht.
\end{definition}

\section{Bewertung und Engineering Decision}

Die Bewertung fragt, ob Kandidat und Evidenz für einen Zweck unter
identifizierten Regeln, Lifecycle und Gültigkeit anwendbar sind. Eine
Engineering Decision kann danach Auswahl oder Ablehnung mit Autorität
dokumentieren. Solver, AXTRAN-Implementierung, App und Thesis besitzen diese
Autorität nicht.

Damit lautet die vollständige Folge
\[
\begin{aligned}
S_{\mathrm{rail}}\to s(t)\to M\to x_M^{\mathrm{pred}}
&\to\text{Vergleich}\to J,\mathcal F\to C\\
&\to\text{Bewertung}\to\text{Engineering Decision}.
\end{aligned}
\]
Geometrie, Modell, Ziel, Kandidat und Entscheidung bleiben getrennt
verantwortet.
